% paper_text.tex -- GENERATED, do not edit by hand.
%
% Draft replacement text for the 21-cm signal-loss and antenna-position
% results, in manuscript order. Regenerate by running, in this order,
%   horizon_position/notebooks/horizon_shift.ipynb
%   horizon_position/notebooks/signal_loss.ipynb
%   horizon_position/notebooks/beam_comparison.ipynb
% in the mock_analysis repo. signal_loss publishes its substitution values
% into signal_loss.npz; beam_comparison runs last, merges them with its own
% and substitutes this template. EVERY number here is computed at generation
% time from the arrays the figures are drawn from, so re-running after the
% 21-cm ensemble changes updates the prose and the figures together and they
% cannot drift apart. Keep it that way -- the one block that ever quoted
% hand-entered numbers (the rotation paragraph, removed 2026-08-25) was also
% the one whose numbers went stale.
%
% This file replaced signal_loss_text.tex.in and beam_comparison_text.tex.in
% on 2026-08-26, when beam_comparison.pdf took over as Fig. 1. Two templates
% substituted by two notebooks could not put the blocks in reading order, and
% the old single-panel caption survived in a file that still looked
% authoritative -- which is how a stale caption gets pasted.
%
% ===================================================================
% FIG. 1 IS NOW beam_comparison.pdf, three panels, under the existing
% \label{fig:singular_values}. The old single-column signal_loss.pdf is
% retired: its content is the CENTRAL PANEL of the new figure, reproduced
% from the same simulation to better than a per cent, so printing both would
% have printed one panel twice. signal_loss.ipynb still runs -- it derives
% N_ANCHOR and every number in blocks 1, 3, 4 and 5 -- it just no longer
% contributes a figure. Do not delete that notebook.
%
% The figure is 7in wide and needs `figure*'. That is the one real cost
% of the replacement: the paper loses a single-column figure early on.
%
% Block 1 is one continuous passage in "Minimising Covariance with the 21-cm
% Signal": the three beams are introduced before the figure is referenced, the
% numbers are attributed to the antenna they belong to as they are quoted, and
% the caveat paragraph closes the section once. It was two blocks until
% 2026-08-26; do not split it apart again.
%
% ===================================================================
% POWER vs AMPLITUDE. Every "per cent of the power" in this file is a ratio of
% squares -- LEADPCT, SPIKEPCT, the LEAK fractions, the DEC decades and the
% TROUGH shares all are. Every comparison of one curve against another is a
% ratio of RMS and must say RMS or amplitude, never power: the figures' y-axis
% is an RMS, the compressions COMPFG/COMPU are RMS ratios, and SPECM is
% 0.1 x (median RMS)/(systematic RMS) x TOPMAG, so "a tenth" there is a tenth
% in amplitude. Calling that a tenth of the POWER would put the position
% requirement at ~0.3 m instead of SPECM m. RETISO/RETBOW/RETVIV are RMS
% fractions too. Check which one you mean before writing either word.
%
% The vocabulary is the manuscript's own: the section has always described
% this as filtering foreground eigenmodes, and every figure labels its axes
% "Foreground modes filtered" and "Residual RMS [K]", so prose, captions and
% axes say the same thing. Block 1 defines $N$ as the number of modes filtered
% and every later block uses that notation. Where the geometry is the argument
% rather than the vocabulary -- the own-basis floors of block 4 -- the words
% "basis" and "subspace" are load-bearing and should stay.
%
% The 10 mK framing is deliberately gone. The manuscript used to say "to filter
% foregrounds to 10 mK, we need to remove the first eight modes", under a
% figure that drew a 10 mK line, and that is what the referee read as a
% sensitivity claim. The residual at that depth is still quoted, as the
% dynamic range the foregrounds are described to; it is no longer a target
% reached, and no figure draws a threshold.
%
% NO FIGURE MARKS A DEPTH: no vertical line at $N$ = NA, no colour scale keyed
% to it, no frequency-space panel drawn at one. That is deliberate, and it is
% also why the three-panel Fig. 1 does not reintroduce the problem the old
% three-panel version had -- those panels were three DEPTHS with a labelled N,
% these are three ANTENNAS with none. Blocks 1 and 4 quote numbers at
% $N$ = NA and NHAND, but a caption or sentence must never imply a figure
% singles either out.
%
% TWO DEPTHS, AND THEY ARE NOT THE SAME ONE. N = NA is where the foreground
% residual drops below the median retained 21-cm signal and stays there -- a
% foreground-vs-SIGNAL crossing, and block 1's headline. N = NHAND, one mode
% inside it, is where an unmodelled TOPMAG m vertical error overtakes the
% foreground floor -- a foreground-vs-SYSTEMATIC crossing, and the depth
% block 4 quotes every position number at. They coincided until the 21-cm
% ensemble was put into antenna temperature; that moved NA up by one and left
% NHAND alone, since neither curve in it involves the 21-cm signal. Do not
% collapse them back into one number.
%
% Block 4 quotes at NHAND because that is where the position error is the term
% limiting the residual. It is also the only place the numbers mean anything:
% SPIKEPCT per cent of what a TOPMAG m vertical error leaves after NHAND modes
% is the single SPIKEth mode, so one more filtered mode drops it by a factor of
% ten. Floors read at NA describe the remainder after that mode, not the error.
% That fragility is block 4's closing argument rather than an embarrassment --
% a fixed-depth filter set either side of one eigenmode reports a different
% result, which is the case for forward modelling.
%
% THE 21-CM ENSEMBLE IS IN ANTENNA TEMPERATURE. Every figure works in
% uncorrected antenna temperature -- ground pickup in, receiver out -- so the
% models are multiplied by the beam-weighted open-sky fraction
% eta = 1 - f_gnd before they are filtered, using EACH ANTENNA'S OWN eta in
% Fig. 1. An isotropic signal reaches that observable attenuated; comparing a
% sky-referred T21 against an antenna-temperature residual would overstate the
% retained signal by 2.2x and was the state of this text before 2026-08-26.
% MED and MEDHAND are therefore NOT sky-referred amplitudes, and any sentence
% that puts them beside an EDGES/SARAS/REACH number must say so. Ground-loss
% correcting instead was considered and rejected: it needs the beam knowledge
% block 1 claims not to use, it makes the foregrounds less compressible, and
% it gives the same answer anyway.
%
% THE VIVALDI IS NOT A REJECTED CANDIDATE, and block 1 lives or dies on this.
% It is a HERA Phase II feed (\citealt{2021ITAP...69.8143F},
% \citealt{2024PASP..136d5002B}), it is what EIGSEP uses for the ground
% antennas (section~\ref{sec:vivaldis}), and it is the antenna the October
% 2024 suspension flew -- the SECOND of the three (July 2024 prototype
% platform, October 2024 Vivaldi, July 2025 bowtie), not the first. So block 1
% frames the comparison as a documented evolution -- a prototype design that
% was flown and was never optimised for low chromaticity, against a bowtie
% developed afterwards against that criterion -- and NOT as a choice between
% two candidates. The
% framing is HISTORICAL, not exculpatory: an earlier revision defended the
% Vivaldi in its own right ("where low chromaticity is not required", "the
% ground antennas have a different job") and that reads as a talk aside, and
% as protesting too much. Saying what each antenna was designed for does the
% same work in fewer words. Keep "used in isolation, without a dish": it is
% the manuscript's own framing in section~\ref{sec:vivaldis}, which the reader
% has not reached yet at this point in the paper, and the word "optimised"
% attached to the bowtie is what carries the contrast.
%
% THE APERTURE POINT NOW ARRIVES AS BOOKKEEPING, not as an argument. ETAVIV >
% ETABOW is asserted in beam_comparison.ipynb, and the fact is stated in the
% units paragraph -- the Vivaldi feed carries proportionally more 21-cm
% amplitude into its filter -- where it is needed anyway to explain why the
% ensemble is attenuated at all. It used to be a standalone sentence ("begins
% with more signal and still ends with less"), which was the most persuasive
% line in the block and also the one in a voice the rest of the paper does not
% use. Stating the eta values BEFORE the retained fractions does the same job:
% the reader has the disadvantage in hand when the fractions land. What this
% must never say is that the Vivaldi ends with less in ABSOLUTE retained
% amplitude -- eta divides out of the fractions, and in absolute terms the two
% antennas are nearly tied.
%
% RETAINED FRACTION AT MATCHED SUPPRESSION, and both halves are load-bearing.
% Fraction, because absolute retained amplitude mixes in aperture efficiency
% and the two antennas are then nearly tied -- a statement about collecting
% area, not spectral overlap. Matched suppression, because at a fixed MODE
% COUNT the Vivaldi appears to retain MORE, for the trivial reason that it has
% suppressed the foregrounds less.
%
% ROTATIONS. There was a further block, a paragraph quantifying what the drive
% grid costs in model complexity; it was removed on 2026-08-25 because it
% introduced a complication only to argue against it and defer it. Nothing in
% the manuscript overclaims without it. But the section before
% \subsection{Forward Modelling} used to say "project out spatial foreground
% modes by rotating the antenna" and "use this modulation to remove
% foregrounds", which borrows the verbs this file uses for the eigenmode
% filter. Folding rotations into that filter's basis is the one thing we have
% actually tested, and it costs modes and hurts 21-cm retention. Those
% sentences now read "separate" and "constrain". Do not put "project out" or
% "remove" back.
%
% Paste the five blocks into rasti_template.tex as marked. Nothing here is
% \input by the paper -- this file is a staging area, not a dependency, and
% nothing in this repo writes to the paper .tex itself.
% ===================================================================


% ===================================================================
% BLOCK 1 -- section "Minimising Covariance with the 21-cm Signal".
% Replaces the whole passage from "Using a representative beam ..." to
% "... as described in section~\ref{subsec:fwd_modelling}."
%
% This was two blocks until 2026-08-26 -- the filter result, then the antenna
% comparison bolted on after it. That order introduced a three-panel figure and
% then quoted a page of numbers that came only from its central panel without
% ever saying so, and a reader who jumps to the figure at first mention had no
% way to know which panel was being described. The three beams are now
% introduced BEFORE the figure is referenced, and the numbers are attributed as
% they are quoted. Do not split it back apart.
%
% Paragraph order: what was simulated; which three beams and why; what the SVD
% shows, the mechanism behind the ordering, and the dynamic range that answers
% the "one part in 10^4" heuristic; what the 21-cm ensemble retains; caveats.
% Results first, caveats last and once.
%
% THE BLOCK IS NOT SELF-CONTAINED, by design. Two sentences reach back into the
% manuscript paragraphs immediately above it, which block 1 does not replace:
% "comfortably inside the dynamic range set out above" answers the "about one
% part in 10^4" heuristic there, and the units paragraph's "left uncorrected
% here too" points at the 300 K ground sentence in paragraph 1. If either of
% those manuscript paragraphs is rewritten, these two clauses have to move with
% them.
%
% QUOTE THE BOWTIE. It is the antenna that exists, so the millikelvin numbers
% are its; the other two appear as mode counts and ratios. Quoting all three in
% millikelvin triples the numbers for no gain and invites the reader to treat
% the isotropic bound as achievable.
%
% THE CONVOLUTION SENTENCE is the physics and should survive editing. Without
% it the ordering of the three curves is just an observation; with it, it is
% the reason the antenna is a design variable at all. An achromatic beam leaves
% the foregrounds' spectral structure as the sky's alone; a chromatic one mixes
% its own frequency-dependent angular structure in, so the measured spectrum
% belongs to neither in isolation.
%
% KHALICHI ET AL. lives in the CAVEAT paragraph, not the beam list. That
% paragraph is already the "what this figure is not" paragraph, so a broader
% comparison over more geometries arrives as the answer to a limitation just
% admitted, rather than as a citation that deflates the result before the
% reader has seen it. It also closes the door this section must not open: a
% reader asking "what about other antenna designs?" is sent to the paper that
% owns that question instead of to a study we would have to do. Do NOT move it
% back up into paragraph 2, and do not add a second sentence arguing the point.
%
% THE POINTER RUNS FORWARD, and this is a fact about where the sections sit.
% subsec:covariance is at rasti_template.tex:204; subsec:eig_antenna is at 344.
% So this is the FIRST mention of Khalichi et al. in the paper, and it carries
% the full deferral ("a broader eigenmode-based analysis of beam chromaticity,
% foreground complexity and overlap with global 21-cm signal models") plus a
% forward ref to subsec:eig_antenna. The existing sentence at
% rasti_template.tex:354 should therefore SHRINK to a short back-reference --
% "... over the EIGSEP operating band (section~\ref{subsec:covariance}),
% following the analysis of Khalichi et al. 2026 (in prep.)" -- so the same
% 14-word phrase is not printed twice. That edit is in the paper repo, not
% here; if it has not been made, this block duplicates it.
%
% It is a plain-text in-prep citation; there is no bib key.
%
% THE VIVALDI IS NOT A REJECTED CANDIDATE -- see the header. Paragraph 2 now
% does it with design intent rather than defence: "used in isolation, without a
% dish, as for the ground antennas", then one sentence of history -- a
% prototype design, flown in October 2024, never optimised for low
% chromaticity, with the bowtie developed afterwards against that criterion.
% "Optimised" attached to the bowtie is load-bearing: it is the whole contrast,
% and it is why the paragraph no longer needs a clause excusing the Vivaldi.
% Keep the British spelling; the manuscript is consistent (optimised,
% minimising, modelled).
%
% RMS0 IS BEAM-DEPENDENT and the value quoted is the bowtie's: the raw
% waterfall RMS is 691 K for the isotropic beam, 728 K for the bowtie and
% 1138 K for the Vivaldi feed, which is directive enough to see more of the
% bright low-frequency sky and less of the 300 K ground. The sentence sits in
% a paragraph that describes all three beams, so it has to say whose number it
% is. Do not generalise it back.
%
% DYNBOW, NOT DYNRANGE. The dynamic-range sentence is quoted at NBOW, the
% depth the same paragraph just said the bowtie needs, so the whole passage
% sits at one N. It used to be quoted at NAM1 = 9 (as DYNRANGE), which put a
% third mode count into a passage that otherwise uses 10 throughout and read as
% a fourth statistic rather than as the answer to the heuristic above. DYNRANGE
% and FGM1 are still emitted by signal_loss.ipynb and are now UNUSED here, like
% RMS1 and RMS4; that is fine -- the substitution assert only catches
% unsubstituted tokens, not unused keys.
%
% DYNBOW is computed in beam_comparison.ipynb, where NBOW and the per-beam
% residual curves both live, as fg_resid["bowtie"][0] / fg_resid["bowtie"][NBOW].
% Its numerator is the same quantity signal_loss rounds to RMS0, and the two
% notebooks are asserted to agree on the bowtie curve to better than a per
% cent, so the sentence and the RMS0 sentence are consistent.
%
% It is a mantissa-exponent string, not an exponent. The old DYNRANGE used to
% be f"{log10(ratio):.0f}" inside $10^{...}$, which rounded 10^5.6 up to 10^6
% and overstated the dynamic range by a factor of 2.5. Keep _sci().
%
% TWO REFERENCE PLANES IN ONE PARAGRAPH, and the units sentence is what keeps
% them apart. MED, SEP0 and SEP1 are ANTENNA temperature (the ensemble is
% attenuated by eta before filtering); the "deeper than 50 mK" selection behind
% NDEEP is SKY brightness temperature -- signal_loss.ipynb computes it as
% -T21_sky.min(axis=1), because trough depth describes the model, not the
% observation. Both appear in mK within four sentences of each other, three
% paragraphs after the section names EDGES, MIST, REACH and SARAS, so the
% reader arrives with a 500 mK trough in mind. Hence BOTH halves of the units
% sentence: the eta range with "smaller ... by that factor", and the explicit
% "Trough depths are quoted as sky brightness temperature". Cutting either one
% leaves the paragraph quietly mixing the two.
%
% Deliberately absent: the RMS at one and four modes (the steep decline is
% visible in the figure and the dynamic-range sentence carries the claim); the
% retention-vs-trough-depth correlations; the `stays below' rule; and the
% repeated "N is not a filter depth we propose to apply to data", which is
% block 4's argument and is made there, once.
%
% ZEUS21 ATTRIBUTION. Its README asks for one paper unconditionally
% (2023MNRAS.523.2587M), a second only if the UVLF module is used, and a third
% for relative velocities, Lyman-Werner feedback or Pop III stars
% (2025PhRvD.111h3503C). Two of the three is correct here.
% ===================================================================

Using representative beams from our electromagnetic simulations of the EIGSEP
antenna (section~\ref{subsec:eig_antenna}), a model of the horizon profile at
the EIGSEP site (Fig.~\ref{fig:horizon}), and the Global Sky Model
\citep{2017MNRAS.464.3486Z}, we simulated one sidereal day of observation by
EIGSEP over 50--250\,MHz with the forward model described in
section~\ref{subsec:fwd_modelling}.
In the simulations, we assumed a constant ground temperature of 300\,K. This
beam-weighted ground pickup is included in the antenna temperature and not
corrected here.

To separate the effect of the antenna from that of the sky, we repeated the
simulation for three beams observing the same sky behind the same horizon: an
achromatic and isotropic beam, the optimised EIGSEP bowtie, and a Vivaldi feed
used in isolation, without a dish, as for the ground antennas
(section~\ref{sec:vivaldis}).
The Vivaldi feed is a prototype design, flown on the October 2024 EIGSEP
suspension but never optimised for low chromaticity; the bowtie geometry was
developed afterwards against that criterion.

We then computed the singular value decomposition of each antenna temperature,
producing its orthogonal spectral eigenmodes, and used those eigenmodes to
filter the corresponding beam-weighted foregrounds.
In Fig.~\ref{fig:singular_values}, we present the residual RMS after filtering
as a function of the number of eigenmodes filtered, $N$, for each of the three
beams. The first point -- zero modes filtered -- is the RMS of the raw waterfall
pixels about zero, which for the bowtie is approximately 730\,K; because
the decomposition is not centred, the first mode captures the mean antenna
temperature spectrum.
The isotropic beam represents an ideal case with no chromatic pattern coupling
to the foregrounds, and so produces the beam-weighted foregrounds with the
least possible spectral structure: it takes 6 modes to filter them below
1\,mK.
The realistic beam patterns of the EIGSEP bowtie and the Vivaldi feed are
chromatic, and that chromaticity does couple to the foregrounds, increasing the
number of modes required to reach the same level: 10 modes for the bowtie
against 18 for the Vivaldi feed.
This ordering follows from the measurement being a convolution of the beam with
the sky. An achromatic beam leaves the spectral structure of the foregrounds as
the sky's alone, whereas a chromatic beam couples its own frequency-dependent
angular structure into them, so that the antenna temperature carries structure
belonging to neither in isolation and more modes are required to describe it.
Filtering the 10 modes the bowtie requires describes its antenna
temperature to one part in $1.2\times10^{6}$, comfortably inside the dynamic range
set out above.

To assess the spectral overlap between the beam-weighted foregrounds and the
expected cosmological signal, we applied each antenna's filter to an ensemble of
1769 global 21-cm models generated with the Zeus21
code\footnote{\url{https://github.com/ZeusCosmo/Zeus21}}
\citep{2023MNRAS.523.2587M, 2025PhRvD.111h3503C}.
The beam-weighted ground pickup is left uncorrected here too, and that applies
to the models as well as to the foregrounds: an isotropic signal reaches the
antenna temperature through the fraction of the beam that sees open sky, which
for the bowtie is 0.36--0.55 across the band, so the retained
amplitudes quoted below are smaller than the corresponding sky brightness
temperatures by that factor. Trough depths are quoted as sky brightness
temperature.
That fraction depends on the beam: averaged over the band it is 0.70 for
the Vivaldi feed against 0.45 for the bowtie, so the Vivaldi feed carries
proportionally more 21-cm signal amplitude into its filter.
For the bowtie, the filter explains foreground variance more effectively than
the cosmological signal: after $N=10$ modes are removed, the RMS of the
median retained 21-cm signal is larger (0.87\,mK) than the RMS of the
residual foregrounds (0.62\,mK).
Compared at the same level of foreground suppression, the bowtie retains
8 per cent of the unfiltered 21-cm RMS, against 4 per cent for
the Vivaldi feed and 13 per cent for the isotropic beam --- that is,
2.2 times the signal the Vivaldi feed retains and 58 per cent
of what an achromatic antenna would.
What the filter leaves also remains informative: among the 1547 models with
absorption troughs deeper than 50\,mK, filtering reduces the median separation
between a pair of models over 70--130\,MHz from 27 to 1.8\,mK, but
78 per cent of pairs stay separated by more than 1\,mK.

While Fig.~\ref{fig:singular_values} is a useful heuristic probing the
covariance between the beam-weighted foregrounds and the 21-cm signal, it does
not represent the complete analysis needed for foreground separation.
Unmodelled instrumental systematics may require more modes filtered than in
this figure.
On the other hand, knowledge of the sky or the beam -- e.g. through beam
mapping (section~\ref{subsec:beam_mapping}) -- may allow chromatic structure in
the beam to be modelled rather than filtered.
A broader eigenmode-based analysis of beam chromaticity, foreground complexity
and overlap with global 21-cm signal models, spanning the range of antenna
geometries considered for EIGSEP, will be presented in Khalichi et al. 2026 (in
prep.); section~\ref{subsec:eig_antenna} describes the geometry it selected.
In practice, EIGSEP will use a differentiable forward model to evaluate
systematics and separate the foregrounds, as described in
section~\ref{subsec:fwd_modelling}.


% ===================================================================
% BLOCK 2 -- caption for beam_comparison.pdf, which is now Fig. 1. Keep
% \label{fig:singular_values}: blocks 2, 4, 5 and the horizon_shift caption all
% refer to it. Two-column `figure*'.
%
% The caption MUST say the grey bands differ between panels. Each antenna has
% its own open-sky fraction and its own eigenbasis, so the same models project
% differently; a reader who assumes it is one band drawn three times makes
% exactly the error a single-panel version would have forced. It must also say
% the isotropic residual leaves the frame, or the truncated curve reads as
% missing data.
%
% What it must NOT do is explain WHY the bands differ beyond naming the cause.
% The three-panel layout carries the point; a sentence arguing it invites the
% argument.
%
% NOT SAID ANY MORE, ANYWHERE: that retained RMS is not retained signal SHAPE.
% An early version showed this in frequency space; that panel is long gone. So
% "median retained signal 0.87 mK" is available to be misread as a surviving
% absorption trough. It is not one: TROUGHNA per cent of the surviving power
% sits inside 70--130 MHz against TROUGHFLAT per cent for a residual with no
% frequency preference at all -- barely a preference -- the rest being band-edge
% ringing from a truncated smooth basis. One clause in either place fixes it.
% Flagged 2026-08-25, left as the author's call.
% ===================================================================

\caption{Residual RMS of the simulated antenna temperature (black) and of the
1769 global 21-cm models filtered with the same modes (5--95 per cent of
ensemble in grey band, median dashed), as a function of the number of
eigenmodes filtered ($N$), for three antennas observing the same sky behind the
same horizon. \textit{Left:} an isotropic beam, which has no chromaticity and
so isolates what the sky alone contributes. \textit{Centre:} the EIGSEP bowtie.
\textit{Right:} a Vivaldi feed of the type used for the ground antennas,
without the dish it was designed to illuminate. Each antenna is filtered in its
own eigenbasis and its models are scaled by its own beam-weighted open-sky
fraction, so the grey bands differ between panels. The isotropic residual falls
below the axis near $N=6$, where it reaches machine precision.}


% ===================================================================
% BLOCK 3 -- caption for horizon_shift.pdf, replacing the existing one.
% Keep \label{fig:horizon_shift} and the figure* environment.
%
% Short and descriptive by design: what the panels show, then the comparison a
% reader who never reaches the text needs. The nuance is block 4's job. Do not
% append a "this is not a proposed analysis" sentence -- blocks 2 and 5 both
% make the point, and the sentence read as an apology. The residual row carries
% the same axes as Fig.~\ref{fig:singular_values} and its grey band is
% described exactly as block 3 describes it -- same ensemble, same artist.
%
% MODE COUNTS HERE, NOT MILLIKELVIN, and this is a reversal. The caption used
% to quote the floor a displacement induces against the median retained signal.
% Printed a page from "the foregrounds leave FG mK", that invites the reading
% that a TOPMAG m position error is a bigger effect than the entire foreground
% sky. It is not: the basis is the foregrounds' own right singular vectors, so
% they compress optimally by construction, and a displaced sky decomposed in
% ITS own basis comes back at OWNLO-OWNHI mK. That argument is block 4,
% paragraph 2, and a caption cannot carry it.
%
% The counts are CLEAR_E / CLEAR_N / CLEAR_U, worst-LST against the median.
% Say "every LST": pooled over LST the same crossings are different, and the
% caption's numbers are not reproducible without that word.
%
% Do NOT write "the change is foreground-like". That is the old manuscript's
% phrase, and block 4 opens by refuting it. The vertical shift needs CLEAR_U
% modes to fall under the median signal where the whole foreground sky needs
% NA, so it is marginally HARDER to filter than the foregrounds, not more
% foreground-like. East and north genuinely are.
%
% CLEAR_U and SPIKE are computed independently -- a stays-below crossing and
% the index of the dominant leftover mode -- and happen to be equal, which is
% WHY the crossing lands where it does and what the "where" clause leans on.
% signal_loss.ipynb asserts they still coincide; if that assert fires, this
% clause needs rewriting rather than re-substituting.
%
% CUT 2026-09-01 BY THE AUTHOR: the trailing "where SPIKEPCT per cent of the
% power surviving NHAND modes lies in the SPIKEth mode alone" clause is gone
% from the manuscript, together with block 4's closing paragraph, which was
% the only text it supported. The caption now ends at "... and CLEAR_U for the
% vertical shift." Do not paste the clause back on its own -- a statistic with
% no argument behind it is worse than neither. The CLEAR_U/SPIKE assert in
% signal_loss.ipynb now gates no published sentence; keep it as a sanity check.
% ===================================================================

\caption{Change in the simulated antenna temperature,
$\Delta T_{\text{ant}}$, when the suspended antenna is displaced to the east
(left column), north (middle column), and up (right column). The displacement
shifts the local horizon, shown in Fig.~\ref{fig:horizon}(b), thereby changing
the fraction of sky occulted by the canyon walls and hence $T_{\text{ant}}$.
\textit{Top row:} $\Delta T_{\text{ant}}$ for a $+1$\,m displacement.
Each curve corresponds to a different LST (one per sidereal hour), coloured as
indicated by the colour bar. \textit{Bottom row:} The RMS over frequency of
$\Delta T_{\text{ant}}$ after filtering the leading eigenmodes of the
unperturbed antenna temperature -- the same modes as the central panel of
Fig.~\ref{fig:singular_values} -- for displacements of 0.1, 1
and 10\,m, coloured light to dark with increasing displacement
magnitude, with all 24 LSTs drawn for each. Grey shows the 21-cm models
of Fig.~\ref{fig:singular_values} retained under the same filter (5--95 per
cent of ensemble, median dashed).
A 1\,m displacement changes $T_{\text{ant}}$ by up to 9.5\,K,
but the change lies largely in the leading eigenmodes: filtering brings every
LST below the median retained 21-cm signal after 7 modes for the
eastward displacement, 4 for the northward, and 10 for the
vertical shift, where 99 per cent of the power surviving
9 modes lies in the 10th mode alone.}


% ===================================================================
% BLOCK 4 -- section "Forward Modelling". Insert after "The tension in the
% rope supporting the antenna limits vertical motion and, due to its
% orientation, east/west motion." -- AFTER that sentence, not before it, since
% paragraph 4 builds directly on what the rope does and does not constrain.
%
% Five short paragraphs, results first: (1) how much a displacement changes
% T_ant and what the filter leaves; (2) why the vertical differs in KIND from
% the horizontals; (3) how direction-dependent the cost is, and the lucky
% alignment with what the suspension actually constrains; (4) the consequence
% for the forward model's prior -- the point of the whole block; (5) why a
% fixed-depth residual is not a detection.
%
% Paragraph 4 is the destination. Everything before it exists to earn the
% asymmetric prior: wide in the horizontal, narrow in the vertical. If the
% block is cut further, cut towards that sentence. Note the framing it turns
% on: what the figure shows is structure left by ignorance, not by an error
% term already in the model. A displacement that is KNOWN is part of the
% forward model and costs nothing, so the quantity that matters is the width
% of the prior, not the size of the displacement.
%
% Paragraph 1 says the modes "carry no knowledge of the displacement" because
% that one clause does the work the old basis-mismatch paragraph did at
% length: it stops the floors being read as "a 1 m position error is worse
% than the entire foreground sky". The own-basis numbers that proved it
% (OWNLO/OWNHI, CVPEN) are still computed and printed by the notebook.
%
% NO position-monitoring requirement is adopted here. An earlier draft said
% the vertical position must be known to SPECM m; that reads as a metrology
% spec, when the real statement is about how much vertical uncertainty the fit
% can carry. Paragraph 4 gives the ladder (MAGLO / TOPMAG / MAGHI m) and lets
% the reader draw the line. SPECM is still computed if it is ever wanted.
%
% Do NOT write "a lesson the field has already drawn from the scrutiny of
% published detections" before the three citations in paragraph 5, or any
% variant. It reads as a dig at EDGES. The citations carry the point.
%
% CUT 2026-09-01 BY THE AUTHOR: paragraph 5 -- "This is why we forward-model
% the instrument ..." through the three citations -- is gone from the
% manuscript. It read as pedagogy rather than result, and the positive case
% for forward modelling is already carried by paragraph 4 (known displacement
% absorbed, unknown marginalised, sensitivities set the prior width). The
% citations survive at rasti_template.tex:161 and :236, so nothing is lost
% bibliographically. Block 3's caption clause was cut with it; restoring
% either one alone is wrong.
% ===================================================================

A 1\,m displacement changes the antenna temperature by up to
5.2\,K to the east, 0.9\,K to the north and 9.5\,K
upward, in each case peaking at the bottom of the band. Under the filter of
Fig.~\ref{fig:singular_values}, whose modes come from the unperturbed antenna
temperature and carry no knowledge of the displacement, an unmodelled
1\,m error raises the residual floor at $N=9$ --- the depth
at which the vertical error first overtakes the foregrounds as the term
limiting the residual --- from 1.82\,mK to 3.52\,mK for a
vertical shift and to 2.04\,mK for an eastward one, and leaves it
unchanged at 1.81\,mK for a northward one. The median retained 21-cm
signal is 1.16\,mK there: the vertical error alone leaves a residual
several times that, the eastward error one comparable to it, and the northward
error one well below it.

The vertical direction differs in kind from the other two. A vertical shift
lowers the local horizon at every azimuth at once (Fig.~\ref{fig:horizon}b),
so its effect on the occulted sky fraction accumulates instead of cancelling,
and the residual it leaves grows roughly in proportion to the displacement and
is symmetric in its sign. A horizontal shift moves the antenna towards some
canyon walls and away from others, so the horizon rises at some azimuths and
falls at others and the two largely cancel, leaving a response that depends on
where the cliff edges fall in azimuth and is neither proportional to the
displacement nor symmetric in its sign.

The cost of a position error is therefore strongly direction-dependent. At
10\,m the induced floor is 30\,mK vertically against
11\,mK eastward and 2.4\,mK northward, so a 10\,m
error in the northward position costs less than a 1\,m error in the
vertical. The alignment with the suspension is fortunate: the rope tension
constrains the vertical and east/west directions and leaves north/south the
least constrained, and the platform IMUs (section~\ref{subsec:box}) cannot
sense lateral motion at all --- but north/south is precisely the direction in
which position error is nearly free.

What Fig.~\ref{fig:horizon_shift} shows is the structure left behind by a
displacement the model does not know about. A displacement that is known is
part of the forward model and is absorbed by the fit; what leaves structure is
the ignorance. The floors quoted above therefore measure modelling error and
not a harder instrument: decomposed in a basis built from itself, each
displaced sky returns to the unperturbed floor --- 1.66--2.03\,mK
even at 10\,m, against 1.82\,mK unperturbed --- so none of the
excess is intrinsic to observing from the displaced position, and the floors
are the bound set by assuming no knowledge of the displacement at all. That
ignorance is marginalised over, and these sensitivities set how wide the prior
may be. Horizontally it may be wide at almost no cost.
Vertically it may not: a 0.1\,m uncertainty leaves the floor at
1.84\,mK, 1 per cent above the unperturbed value, but by
1\,m the floor is 3.52\,mK, three times the median retained
signal, and by 10\,m it is 30\,mK, more than an order of
magnitude above it.

This is why we forward-model the instrument rather than filtering to a fixed
depth and interpreting what is left. An unmodelled displacement of the size we
must anticipate leaves a residual of the same order as the retained signal and
similar to it in shape: at $N=9$ the leftover from a 1\,m
vertical error is 2.6 times the median 21-cm model in RMS, and their
cosine similarity reaches 0.98 (median 0.63). Nor is the choice of
depth a way out. Almost all of what the vertical error leaves at that depth is
a single eigenmode --- 99 per cent of it, the 10th, worth
3.18\,mK --- so filtering one further mode drops the leftover to
0.36 times the median retained signal, and a filter set either side of
that one mode reports a qualitatively different result. An excess with respect
to a foreground model is only as trustworthy as the instrument model behind it
\citep{2018Natur.564E..32H, 2019ApJ...874..153B, 2019ApJ...880...26S}.


% ===================================================================
% BLOCK 5 -- section "Forward Modelling". Replaces the sentence beginning
% "To quantify the effect of the horizon coupling, we computed the change in
% antenna temperature ... by 1 m to the east, north, and up ...".
%
% Needed because Fig.~\ref{fig:horizon_shift}'s residual row now shows three
% displacement magnitudes, and only 1 m was previously introduced. No new
% simulations: run_sims.py has always produced +/-0.1, 1 and 10 m for each
% axis (see positions.py), and the other magnitudes were simply unplotted.
% ===================================================================

To quantify the effect of the horizon coupling, we computed the change in
antenna temperature, $\Delta T_{\text{ant}}(\nu)$, caused by displacing the
antenna to the east, north, and up, relative to the nominal position used for
the horizon profile of Fig.~\ref{fig:horizon}. Each axis was displaced by
0.1, 1 and 10\,m, spanning the range from a
displacement we expect to be able to monitor to one large enough to be
obviously unacceptable, so that the dependence on displacement can be read
off rather than inferred from a single case.
